\documentclass[11pt,a4paper]{article}
\usepackage{turkos}

\begin{document}
\phaseheader{3}{Segmentation, the IDT and CPU Exceptions}{Install your own descriptor tables and turn every CPU fault into a readable crash report instead of a silent reboot}{16}{22}

\ataglance{2 weeks}{4}{Phase 2: C kernel with \code{kprintf}, serial logging and \code{panic}}{Your own Global Descriptor Table
(flat model with kernel and user segments), an Interrupt Descriptor Table with 32 exception stubs in NASM, a common C
handler that receives the full register state, and a crash screen that names the exception and dumps the registers.}

\section{Why this phase}

Right now, if your kernel divides by zero or jumps to a bad address, the CPU looks for an exception handler, finds
none, faults again, and resets the machine. You see QEMU reboot and nothing else. After this phase the same bug
prints ``Divide error (\#DE) at EIP=0x00101234'' with every register, and you can go straight to the faulty line.
No other change in the whole project saves as much debugging time.

The phase has two parts that belong together. The \textbf{GDT} defines the segments the CPU uses and, more
importantly for us, the privilege levels: ring 0 for the kernel and ring 3 for users. The Multiboot contract says
GRUB's GDT may be invalid, so you have to install your own before touching a segment register. The \textbf{IDT} tells
the CPU where to jump for each of the 256 interrupt vectors. Vectors 0--31 are CPU exceptions, which you handle now.
Vectors 32--47 will be hardware interrupts in Phase 4, and vector \code{0x80} becomes the system call gate in
Phase 10. This is the \emph{trap table} that \OSTEP chapter 6 describes: the kernel installs it at boot, and from
then on the hardware knows where to go when something needs the kernel's attention.

\section{Concepts you need}

\subsection{Segmentation, and why we switch it off}

In protected mode every memory access goes through a segment. A \textbf{segment selector} in CS, DS, SS and the
other segment registers picks a \textbf{descriptor} from the GDT (Figure~\ref{fig:selector}), and the descriptor
gives the segment's base address, its limit, and its access rights. The CPU adds the base to every address and checks
it against the limit. \OSTEP chapter 16 explains the general idea, and \OSDI\,\S4.6.2 describes the Intel version.

\begin{figure}[htbp]
\centering
\begin{tikzpicture}[x=0.55cm, y=1cm]
  \draw[draw=tkNavy!70, fill=tkBlue!10] (0,0) rectangle (13,0.7);
  \draw[draw=tkNavy!70, fill=tkAmber!12] (13,0) rectangle (14,0.7);
  \draw[draw=tkNavy!70, fill=tkRed!10] (14,0) rectangle (16,0.7);
  \node[font=\sffamily\scriptsize] at (6.5,0.35) {index into the descriptor table};
  \node[font=\sffamily\scriptsize] at (13.5,0.35) {TI};
  \node[font=\sffamily\scriptsize] at (15,0.35) {RPL};
  \foreach \x/\b in {0.5/15, 12.5/3, 13.5/2, 14.5/1, 15.5/0}
    \node[font=\sffamily\scriptsize, text=tkGray] at (\x,0.95) {\b};
  \node[tknote, text width=8.5cm, align=left, anchor=north west] at (0,-0.15)
    {TI = 0 selects the GDT. RPL = requested privilege level. Selector \code{0x08} = index 1, ring 0 (kernel code); \code{0x10} = index 2 (kernel data); \code{0x1B} = index 3 with RPL 3 (user code).};
\end{tikzpicture}
\caption{A 16-bit segment selector.}
\label{fig:selector}
\end{figure}

Modern operating systems use paging for memory protection, so they neutralise segmentation with the \textbf{flat
model}: every descriptor has base 0 and a limit of 4\,GiB, which makes logical addresses equal to linear addresses.
The GDT still matters for three reasons. The CPU refuses to run in protected mode without one. The \textbf{privilege
level} of running code is the DPL of its code segment, so ring 0 and ring 3 need separate descriptors. And the
\textbf{TSS} descriptor that Phase 10 needs for user mode lives in the GDT.

Each descriptor packs its fields into 8 bytes in a famously awkward layout (Figure~\ref{fig:desc}), a leftover from
the 286. The two bytes you choose are the \textbf{access byte} (present bit, DPL, code or data, readable or writable)
and the \textbf{flags} nibble (4\,KiB granularity and 32-bit operation). Table~\ref{tab:gdt} lists the values
Turk-OS uses.

\begin{figure}[htbp]
\centering
\begin{tikzpicture}[x=0.44cm, y=0.8cm]
  % high dword
  \draw[draw=tkNavy!70, fill=tkGray!10] (0,1) rectangle (8,1.8);
  \draw[draw=tkNavy!70, fill=tkGreen!10] (8,1) rectangle (12,1.8);
  \draw[draw=tkNavy!70, fill=tkGray!10] (12,1) rectangle (16,1.8);
  \draw[draw=tkNavy!70, fill=tkRed!10] (16,1) rectangle (24,1.8);
  \draw[draw=tkNavy!70, fill=tkGray!10] (24,1) rectangle (32,1.8);
  \node[font=\sffamily\scriptsize] at (4,1.4) {base 31:24};
  \node[font=\sffamily\scriptsize] at (10,1.4) {G D/B L A};
  \node[font=\sffamily\scriptsize] at (14,1.4) {limit 19:16};
  \node[font=\sffamily\scriptsize] at (20,1.4) {P DPL S type};
  \node[font=\sffamily\scriptsize] at (28,1.4) {base 23:16};
  % low dword
  \draw[draw=tkNavy!70, fill=tkGray!10] (0,0) rectangle (16,0.8);
  \draw[draw=tkNavy!70, fill=tkGray!10] (16,0) rectangle (32,0.8);
  \node[font=\sffamily\scriptsize] at (8,0.4) {base 15:0};
  \node[font=\sffamily\scriptsize] at (24,0.4) {limit 15:0};
  \node[font=\sffamily\scriptsize\bfseries, anchor=east] at (-0.3,1.4) {bytes 4--7};
  \node[font=\sffamily\scriptsize\bfseries, anchor=east] at (-0.3,0.4) {bytes 0--3};
  \node[tknote] at (10,2.1) {flags};
  \node[tknote] at (20,2.1) {access byte};
  \foreach \x/\b in {0.4/31, 7.6/24, 15.6/16, 16.4/15, 23.6/8, 31.6/0}
    \node[font=\sffamily\tiny, text=tkGray] at (\x,-0.25) {\b};
\end{tikzpicture}
\caption{The 8-byte segment descriptor. Base and limit are scattered across both halves.}
\label{fig:desc}
\end{figure}

\begin{table}[htbp]
\centering\small\renewcommand{\arraystretch}{1.25}
\begin{tabular}{@{}c l l l c c@{}}
\toprule
\tkth{Index} & \tkth{Selector} & \tkth{Purpose} & \tkth{Access byte} & \tkth{Flags} & \tkth{Limit} \\
\midrule
0 & \code{0x00} & null descriptor (required, never used) & \code{0x00} & \code{0x0} & 0 \\
1 & \code{0x08} & kernel code: present, ring 0, executable, readable & \code{0x9A} & \code{0xC} & \code{0xFFFFF} \\
2 & \code{0x10} & kernel data: present, ring 0, writable & \code{0x92} & \code{0xC} & \code{0xFFFFF} \\
3 & \code{0x18} & user code: as kernel code but DPL 3 (Phase 10) & \code{0xFA} & \code{0xC} & \code{0xFFFFF} \\
4 & \code{0x20} & user data: as kernel data but DPL 3 (Phase 10) & \code{0xF2} & \code{0xC} & \code{0xFFFFF} \\
5 & \code{0x28} & task state segment (filled in Phase 10) & \code{0x89} & \code{0x0} & \code{sizeof(tss)-1} \\
\bottomrule
\end{tabular}
\caption{The Turk-OS GDT. Flags \code{0xC} = 4\,KiB granularity + 32-bit, so limit \code{0xFFFFF} covers 4\,GiB.}
\label{tab:gdt}
\end{table}

\subsection{Interrupts, exceptions and traps}

Three kinds of event make the CPU stop and run a handler. \textbf{Hardware interrupts} come from devices (Phase 4).
\textbf{Exceptions} come from the CPU itself when an instruction can't complete. \textbf{Software interrupts} come
from the \code{int n} instruction, which is how user programs will call the kernel. Exceptions are further divided
by what happens on return (Table~\ref{tab:exc}). After a \textbf{fault}, the saved EIP points at the faulting
instruction, so returning retries it; that is how page faults will work in Phase 6. After a \textbf{trap}, EIP points
past the instruction, so execution continues; \code{int3} works this way. An \textbf{abort} means the machine state is
unreliable, and all you can do is report it and stop.

\begin{table}[htbp]
\centering\small\renewcommand{\arraystretch}{1.2}
\begin{tabular}{@{}r l l l c@{}}
\toprule
\tkth{Vector} & \tkth{Mnemonic} & \tkth{Name} & \tkth{Class} & \tkth{Error code} \\
\midrule
0 & \#DE & Divide error & fault & no \\
1 & \#DB & Debug & fault or trap & no \\
2 & -- & Non-maskable interrupt & interrupt & no \\
3 & \#BP & Breakpoint (\code{int3}) & trap & no \\
4 & \#OF & Overflow (\code{into}) & trap & no \\
5 & \#BR & Bound range exceeded & fault & no \\
6 & \#UD & Invalid opcode & fault & no \\
7 & \#NM & Device (FPU) not available & fault & no \\
8 & \#DF & Double fault & abort & yes (always 0) \\
10 & \#TS & Invalid TSS & fault & yes \\
11 & \#NP & Segment not present & fault & yes \\
12 & \#SS & Stack-segment fault & fault & yes \\
13 & \#GP & General protection & fault & yes \\
14 & \#PF & Page fault (address in CR2) & fault & yes \\
16 & \#MF & x87 floating-point error & fault & no \\
17 & \#AC & Alignment check & fault & yes \\
18 & \#MC & Machine check & abort & no \\
19 & \#XM & SIMD floating-point & fault & no \\
21 & \#CP & Control protection & fault & yes \\
\bottomrule
\end{tabular}
\caption{The CPU exceptions you will meet. Vectors 9, 15, 20 and 22--31 are reserved or rare; give them a generic handler.}
\label{tab:exc}
\end{table}

\subsection{What happens on an interrupt}

When vector $n$ fires, the CPU reads IDT entry $n$. Each entry (a \textbf{gate}) holds the handler's address, the
code segment selector to run it with (\code{0x08}), and a type byte. Turk-OS uses \code{0x8E}: present, DPL 0, 32-bit
\emph{interrupt gate}, which clears the interrupt flag on entry. A \emph{trap gate} (\code{0x8F}) leaves interrupts
enabled. The CPU then pushes EFLAGS, CS and EIP on the stack and, for some exceptions, an error code, and jumps to the
handler. The handler returns with \code{iret}, which pops EIP, CS and EFLAGS back. (When the interrupt arrives in ring
3, the CPU also switches stacks and pushes the user SS and ESP; that is Phase 10's topic.)

C can't be the handler directly: a C function does not know it must preserve every register and return with
\code{iret}. So each vector gets a tiny assembly stub that makes the stack look the same for every interrupt,
saves all registers, and calls one C function with a pointer to the saved state (Figure~\ref{fig:frame}). Getting
this layout exactly right, and matching it with a C struct, is the heart of the phase.

\begin{figure}[htbp]
\centering
\begin{tikzpicture}[slot/.style={draw=tkNavy!70, minimum width=3.6cm, minimum height=0.46cm, font=\ttfamily\scriptsize},
  br/.style={decorate, decoration={brace, amplitude=4pt}, line width=0.6pt}]
  \node[slot, fill=tkGray!12] (ss) at (0,0) {ss, useresp};
  \node[slot, fill=tkRed!10, below=0pt of ss] (ef) {eflags};
  \node[slot, fill=tkRed!10, below=0pt of ef] (cs) {cs};
  \node[slot, fill=tkRed!10, below=0pt of cs] (ip) {eip};
  \node[slot, fill=tkAmber!14, below=0pt of ip] (ec) {err\_code (or dummy 0)};
  \node[slot, fill=tkAmber!14, below=0pt of ec] (in) {int\_no};
  \node[slot, fill=tkBlue!10, below=0pt of in] (ax) {eax ecx edx ebx};
  \node[slot, fill=tkBlue!10, below=0pt of ax] (sp) {esp ebp esi edi};
  \node[slot, fill=tkGreen!10, below=0pt of sp] (ds) {ds};
  \node[slot, fill=white, below=0pt of ds] (pa) {pointer to this frame};
  \draw[br, tkGray] ([xshift=4pt]ss.north east) -- ([xshift=4pt]ss.south east) node[midway, right=5pt, tknote, align=left] {only when coming from ring 3};
  \draw[br, tkRed] ([xshift=4pt]ef.north east) -- ([xshift=4pt]ip.south east) node[midway, right=5pt, tknote, align=left] {pushed by the CPU};
  \draw[br, tkAmber] ([xshift=4pt]ec.north east) -- ([xshift=4pt]in.south east) node[midway, right=5pt, tknote, align=left] {pushed by the per-vector stub\\(the CPU pushes err\_code for 8, 10--14, 17, 21)};
  \draw[br, tkBlue] ([xshift=4pt]ax.north east) -- ([xshift=4pt]sp.south east) node[midway, right=5pt, tknote, align=left] {\code{pusha} in the common stub};
  \draw[br, tkGreen] ([xshift=4pt]ds.north east) -- ([xshift=4pt]ds.south east) node[midway, right=5pt, tknote, align=left] {saved data segment};
  \draw[br, tkNavy] ([xshift=4pt]pa.north east) -- ([xshift=4pt]pa.south east) node[midway, right=5pt, tknote, align=left] {\code{push esp}: the argument \code{struct regs *r}};
  \draw[tkarrow] ([xshift=-12pt]ss.north west) -- ([xshift=-12pt]pa.south west) node[midway, left, tknote, text width=1.6cm] {pushed in this order};
\end{tikzpicture}
\caption{The stack when \code{isr_handler} is called. \code{struct regs} lists these fields from the bottom (\code{ds}) up.}
\label{fig:frame}
\end{figure}

\section{Reading plan}

\begin{readingplan}
\must & \LOB & Ch.~5 \emph{Segmentation}: accessing memory, the GDT, loading the GDT & 33--37 & Tasks 3.1--3.2. \\
\must & \LOB & Ch.~6 \emph{Interrupts and Input}, up to and including \emph{Loading the IDT} & 39--43 & Tasks 3.3--3.5; PIC and keyboard wait for Phase 4. \\
\must & \OSTEP & Ch.~6 \emph{Mechanism: Limited Direct Execution} & 49--62 & Why a trap table exists and what the hardware does on a trap. \\
\must & \OSTEP & Ch.~16 \emph{Segmentation} & 149--160 & The general idea behind the descriptors you are writing. \\
\should & \OSDI & \S4.6.2 \emph{Segmentation with Paging: The Intel Pentium} & 415--420 & Selectors, descriptors, GDT and LDT on real x86 hardware. \\
\should & \MOS & \S3.7.3 \emph{Segmentation with Paging: The Intel x86} & 247--252 & A second explanation, with the history of why x86 kept segments. \\
\should & \MOS & \S5.1.5 \emph{Interrupts Revisited} & 347--351 & Precise and imprecise interrupts; what ``saved state'' means. \\
\should & \OSDI & \S2.6.8 \emph{Interrupt Handling in MINIX 3} & 167--178 & A complete, working version of what you build here. \\
\could & \OSC & \S1.2 \emph{Computer-System Organization} (interrupts); \S17.3 \emph{Protection Rings} & 7--15, 669--671 & Interrupt mechanics and rings in general terms. \\
\end{readingplan}

From the Intel SDM Volume 3A, read the parts of chapter 3 (protected-mode memory management) about segment
selectors and descriptors, and chapter 6 (interrupt and exception handling), especially the exception
classification, the error code format, and the per-exception reference at its end. Keep that last part open while
you write the handler names.

\begin{minixbox}
MINIX builds its GDT and IDT in \texttt{kernel/protect.c (08300)}; look for \code{prot_init}, which fills the tables
from a list of gate descriptions much like your \code{idt_init}. The entry stubs are macros in
\texttt{kernel/mpx386.s (06300)}, and exceptions end up in \texttt{kernel/exception.c (08000)}. There, notice the
decision you will make in Phase 10: an exception in a user process becomes a signal to that process (for example,
divide error sends \code{SIGFPE}), while an exception inside the kernel is a panic.
\end{minixbox}

\section{Build it}

\task{Describe the GDT in C}
Write \code{kernel/arch/i386/gdt.c}. Pack the structs (they are hardware formats), fill the five entries of
Table~\ref{tab:gdt}, and leave room for the TSS.

\begin{lst}{c}{kernel/arch/i386/gdt.c}
struct gdt_entry {
    uint16_t limit_low;
    uint16_t base_low;
    uint8_t  base_mid;
    uint8_t  access;
    uint8_t  gran;          /* high nibble: flags (G, D/B); low nibble: limit 19:16 */
    uint8_t  base_high;
} __attribute__((packed));

struct gdt_ptr {
    uint16_t limit;         /* size of the table in bytes, minus one */
    uint32_t base;          /* linear address of the table           */
} __attribute__((packed));

static struct gdt_entry gdt[6];
static struct gdt_ptr   gp;
extern void gdt_flush(struct gdt_ptr *p);   /* in gdt_flush.s */

static void gdt_set(int i, uint32_t base, uint32_t limit, uint8_t access, uint8_t flags)
{
    gdt[i].base_low  = base & 0xFFFF;
    gdt[i].base_mid  = (base >> 16) & 0xFF;
    gdt[i].base_high = (base >> 24) & 0xFF;
    gdt[i].limit_low = limit & 0xFFFF;
    gdt[i].gran      = ((limit >> 16) & 0x0F) | (flags << 4);
    gdt[i].access    = access;
}

void gdt_init(void)
{
    gdt_set(0, 0, 0,       0x00, 0x0);  /* null                        */
    gdt_set(1, 0, 0xFFFFF, 0x9A, 0xC);  /* 0x08 kernel code            */
    gdt_set(2, 0, 0xFFFFF, 0x92, 0xC);  /* 0x10 kernel data            */
    gdt_set(3, 0, 0xFFFFF, 0xFA, 0xC);  /* 0x18 user code  (Phase 10)  */
    gdt_set(4, 0, 0xFFFFF, 0xF2, 0xC);  /* 0x20 user data  (Phase 10)  */
    /* entry 5 (0x28) becomes the TSS in Phase 10 */
    gp.limit = sizeof(gdt) - 1;
    gp.base  = (uint32_t)&gdt;
    gdt_flush(&gp);
}
\end{lst}

\task{Load the GDT and reload the segment registers}
\code{lgdt} only tells the CPU where the table is. The segment registers still hold the old descriptors in a hidden
cache until you reload them. Data segments can be loaded with \code{mov}; CS can only be changed by a far jump (or a
far call or \code{iret}).

\begin{lst}{asm}{kernel/arch/i386/gdt_flush.s}
global gdt_flush
gdt_flush:
    mov  eax, [esp + 4]     ; argument: pointer to struct gdt_ptr
    lgdt [eax]              ; load the GDT register
    mov  ax, 0x10           ; 0x10 = kernel data selector
    mov  ds, ax
    mov  es, ax
    mov  fs, ax
    mov  gs, ax
    mov  ss, ax
    jmp  0x08:.reload_cs    ; far jump: CS = 0x08 (kernel code)
.reload_cs:
    ret
\end{lst}
\verify{The kernel keeps running after \code{gdt_init()}. In the QEMU monitor, \code{info registers} shows a
\code{GDT=} line whose base is the address of your \code{gdt} array (check with \code{i686-elf-nm}) and whose limit is
\code{0000002f}; CS is \code{0008} and DS is \code{0010}.}

\task{Build and load the IDT}
Write \code{kernel/arch/i386/idt.c} with a 256-entry table of gates, a setter, and \code{idt_init}, which installs the
32 exception stubs from the next task and loads the table with \code{lidt}. Entries you don't fill stay zero, which
means ``not present''; if such a vector fires you get a \#GP, which your handler will report.

\begin{lst}{c}{kernel/arch/i386/idt.c}
struct idt_entry {
    uint16_t offset_low;    /* handler address bits 15:0            */
    uint16_t selector;      /* code segment to run it in: 0x08      */
    uint8_t  zero;
    uint8_t  type_attr;     /* 0x8E = present, ring 0, 32-bit interrupt gate */
    uint16_t offset_high;   /* handler address bits 31:16           */
} __attribute__((packed));

struct idt_ptr { uint16_t limit; uint32_t base; } __attribute__((packed));

static struct idt_entry idt[256];
extern uint32_t isr_stub_table[32];     /* addresses of isr0..isr31, from isr.s */

void idt_set_gate(uint8_t n, uint32_t handler, uint16_t sel, uint8_t flags)
{
    idt[n].offset_low  = handler & 0xFFFF;
    idt[n].offset_high = handler >> 16;
    idt[n].selector    = sel;
    idt[n].zero        = 0;
    idt[n].type_attr   = flags;
}

void idt_init(void)
{
    for (int i = 0; i < 32; i++)
        idt_set_gate(i, isr_stub_table[i], 0x08, 0x8E);
    struct idt_ptr ip = { sizeof(idt) - 1, (uint32_t)&idt };
    __asm__ volatile ("lidt %0" : : "m"(ip));
}
\end{lst}

\task{Write the interrupt stubs}
This is the most important assembly in the project, so type it yourself and comment every line. Two NASM macros
generate the 32 stubs: one for vectors where the CPU pushes an error code and one that pushes a dummy zero, so every
frame has the same shape. All stubs jump to one common routine that builds the rest of Figure~\ref{fig:frame}.

\begin{lst}{asm}{kernel/arch/i386/isr.s}
extern isr_handler

%macro ISR_NOERR 1
isr%1:
    push dword 0            ; dummy error code: keeps every frame the same shape
    push dword %1           ; the vector number
    jmp  isr_common
%endmacro

%macro ISR_ERR 1
isr%1:
    push dword %1           ; the CPU has already pushed an error code
    jmp  isr_common
%endmacro

ISR_NOERR 0
ISR_NOERR 1
ISR_NOERR 2
ISR_NOERR 3
ISR_NOERR 4
ISR_NOERR 5
ISR_NOERR 6
ISR_NOERR 7
ISR_ERR   8
ISR_NOERR 9
ISR_ERR   10
ISR_ERR   11
ISR_ERR   12
ISR_ERR   13
ISR_ERR   14
ISR_NOERR 15
ISR_NOERR 16
ISR_ERR   17
ISR_NOERR 18
ISR_NOERR 19
ISR_NOERR 20
ISR_ERR   21
ISR_NOERR 22                ; 22..31: reserved, no error code
ISR_NOERR 23
ISR_NOERR 24
ISR_NOERR 25
ISR_NOERR 26
ISR_NOERR 27
ISR_NOERR 28
ISR_NOERR 29
ISR_NOERR 30
ISR_NOERR 31

isr_common:
    pusha                   ; saves eax, ecx, edx, ebx, esp, ebp, esi, edi
    mov  ax, ds
    push eax                ; save the interrupted code's data segment
    mov  ax, 0x10           ; switch to the kernel data segment
    mov  ds, ax
    mov  es, ax
    mov  fs, ax
    mov  gs, ax
    push esp                ; argument for isr_handler: pointer to the frame
    call isr_handler
    add  esp, 4             ; drop the argument
    pop  eax                ; restore the original data segment
    mov  ds, ax
    mov  es, ax
    mov  fs, ax
    mov  gs, ax
    popa
    add  esp, 8             ; drop int_no and err_code
    iret                    ; pop eip, cs, eflags and resume

global isr_stub_table       ; table of stub addresses for idt_init()
isr_stub_table:
%assign i 0
%rep 32
    dd isr%+i
%assign i i+1
%endrep
\end{lst}

\task{The C exception handler}
Mirror the frame in a struct, lowest address first. Then write \code{isr_handler}: if a handler is registered for the
vector (next task), call it and return; otherwise print a crash report and panic. For a page fault, also print CR2,
which holds the faulting address.

\begin{lst}{c}{include/arch/regs.h and kernel/arch/i386/isr.c}
struct regs {
    uint32_t ds;                                        /* pushed last by isr_common */
    uint32_t edi, esi, ebp, esp_at_pusha, ebx, edx, ecx, eax;   /* pusha */
    uint32_t int_no, err_code;                          /* pushed by the stub        */
    uint32_t eip, cs, eflags;                           /* pushed by the CPU         */
    uint32_t useresp, ss;                               /* only from ring 3 (Phase 10) */
};

static const char *exception_names[32] = {
    "Divide error", "Debug", "NMI", "Breakpoint", "Overflow", "Bound range exceeded",
    "Invalid opcode", "Device not available", "Double fault", "Coprocessor overrun",
    "Invalid TSS", "Segment not present", "Stack-segment fault", "General protection",
    "Page fault", "Reserved", "x87 FPU error", "Alignment check", "Machine check",
    "SIMD FP exception", "Virtualization", "Control protection", /* 22..31 */ };

void isr_handler(struct regs *r)
{
    if (handlers[r->int_no]) {              /* registered handler: let it deal with it */
        handlers[r->int_no](r);
        return;
    }
    const char *name = exception_names[r->int_no] ? exception_names[r->int_no] : "Reserved";
    kprintf("\n*** CPU exception %u: %s (error code 0x%x)\n", r->int_no, name, r->err_code);
    kprintf("EIP=%08x  CS=%04x  EFLAGS=%08x\n", r->eip, r->cs, r->eflags);
    kprintf("EAX=%08x  EBX=%08x  ECX=%08x  EDX=%08x\n", r->eax, r->ebx, r->ecx, r->edx);
    kprintf("ESI=%08x  EDI=%08x  EBP=%08x\n", r->esi, r->edi, r->ebp);
    if (r->int_no == 14) {
        uint32_t cr2;
        __asm__ volatile ("mov %%cr2, %0" : "=r"(cr2));
        kprintf("CR2 (faulting address) = %08x\n", cr2);
    }
    panic("unhandled CPU exception");
}
\end{lst}
\verify{\code{sizeof(struct regs)} is 64 bytes (16 fields). Assert it at compile time with
\code{_Static_assert}.}

\begin{tipbox}[Portability note]
Every field of \code{struct regs} is an x86 register. In Phase 16 it becomes \code{struct trapframe}, one per
architecture, and code outside \code{kernel/arch/i386} reads it only through small helpers: the system call number,
an argument, the return value, the program counter. Start the habit now. Leave field names like \code{eax},
\code{eip} and \code{cs} to the x86 code, and when a later phase (page faults, system calls, signals) needs a
register elsewhere, give it a one-line inline helper instead of a field name.
\end{tipbox}

\task{A handler registration API}
Add \code{register_interrupt_handler(uint8_t n, void (*fn)(struct regs *))} backed by a 256-entry table of function
pointers. Phase 4 installs the timer and keyboard handlers through it, Phase 6 the page-fault handler, and Phase 10
the system call handler. Register a breakpoint handler now that prints ``breakpoint at EIP=\dots'' and simply
returns.

\task{Trigger every exception you can}
Add a \code{crash} command to your boot sequence (a real command monitor arrives in Phase 4). Write one tiny
function per exception and check that each gives the expected report. The \code{volatile} stops the compiler from
noticing the division by zero at compile time.

\begin{lst}{c}{tests/exceptions.c}
void test_divide_error(void)  { volatile int zero = 0; volatile int x = 1 / zero; (void)x; }
void test_invalid_opcode(void) { __asm__ volatile ("ud2"); }
void test_breakpoint(void)     { __asm__ volatile ("int3"); kprintf("...and we continued!\n"); }
void test_gp(void)             { __asm__ volatile ("mov $0x2B, %ax; mov %ax, %ds"); } /* entry 5 is still empty */
\end{lst}
\verify{Divide error, invalid opcode and general protection each print a report with a sensible EIP (check it
with \code{i686-elf-addr2line -e build/kernel.elf <eip>}). The breakpoint test prints both messages and the kernel
keeps running.}

\section{Testing and debugging}

Two QEMU tools are essential here. \code{-d int -no-reboot} logs every interrupt and exception with its vector
(\code{v=0d} is \#GP, \code{v=0e} is \#PF) and the error code, and stops instead of rebooting on a triple fault. The
monitor's \code{info registers} shows the GDT and IDT base and limit. If you install Bochs, its debugger can also
decode the tables with \code{info gdt} and \code{info idt}.

\begin{pitfalls}
Reboot right after \code{lgdt} & Wrong limit (should be \code{sizeof(gdt)-1}), struct not packed, or the far jump is missing & Check \code{sizeof(struct gdt_entry) == 8}; look at \code{GDT=} in \code{info registers}. \\
\#GP immediately after loading DS or SS & Selector points at the wrong entry, or the access byte is wrong & \code{0x10} must be index 2 with access \code{0x92}; print the table bytes. \\
The divide-error report prints forever & Returning from a fault re-executes the faulting instruction & Faults must be fixed or must end in \code{panic}; only traps return safely. \\
Crash or garbage after \code{iret} & Unbalanced stack: missing \code{add esp, 8}, or a stub pushed the wrong number of values & Every path must pop exactly what it pushed; step through \code{isr_common} in GDB. \\
\code{int_no} or registers look shuffled & \code{struct regs} field order doesn't match the push order & The struct lists fields from the \emph{last} push to the first. \\
Double fault (vector 8) on every exception & The handler itself faults, for example because of a stack overflow & Use \code{-d int} to see the first exception in the chain; that one is the real bug. \\
\end{pitfalls}

\section{Milestone}

\begin{milestonebox}[No more silent reboots]
\begin{checklist}
  \item \code{info registers} shows your GDT (limit \code{0x2f}) and IDT (limit \code{0x7ff}), with CS=\code{0x08} and DS=\code{0x10}.
  \item Divide by zero, \code{ud2} and a bad segment load each produce a named crash report with registers, and no reboot.
  \item \code{int3} prints a message and execution continues after it.
  \item \code{register_interrupt_handler} works, and \code{sizeof(struct regs)} is checked at compile time.
  \item You can explain every line of \code{isr_common}.
\end{checklist}
\end{milestonebox}

\begin{questionbox}
\begin{enumerate}
  \item If the flat model makes segmentation do nothing, why is a GDT still required, and why does it need ring 3 entries?
  \item Decode selector \code{0x23}: which GDT entry, which table, which RPL?
  \item What does an interrupt gate do to the IF flag that a trap gate does not? Why does Turk-OS want that?
  \item Which exceptions push an error code, and why do the other stubs push a dummy value?
  \item Draw the stack at the moment \code{isr_handler} is called.
  \item Explain the difference between a fault, a trap and an abort, with one example of each.
  \item Why does returning from the divide-error handler run the handler again?
\end{enumerate}
\end{questionbox}

\section{Going further}

Print a backtrace in the crash report by walking the chain of saved EBP values: \code{[ebp]} holds the caller's EBP and
\code{[ebp+4]} the return address. Because you compile with \code{-fno-omit-frame-pointer}, the chain is intact. Print
raw addresses for now and symbolise them on the host with \code{addr2line}; in Phase 15 you can embed a symbol table.
If you want to go deep, give the double fault its own known-good stack through a task gate, so even a kernel stack
overflow produces a report. That trick needs the TSS machinery from Phase 10.

\nextphase{4}{Hardware Interrupts: PIC, Timer and Keyboard}
\booklegend
\end{document}
